\chapter*{How to use this book}
\addcontentsline{toc}{chapter}{How to use this book}
\markboth{How to use this book}{How to use this book}

\section*{What this is}

This is a tutor, not a summary. It teaches one thing: \textbf{how a neural chess engine
thinks, from the ground up, in enough detail that you could rebuild it.} By the end you
should be able to look at a line of Leela's verbose output --- something like
\begin{quote}\small\ttfamily
info string b3b8 (N: 214) (P: 3.11\%) (Q: 0.99871) (U: 0.04102)
\end{quote}
and read it the way you read a chess diagram: instantly, structurally, knowing what every
number is, where it came from, what would change it, and what it would look like if it
were lying to you.

We get there by building. Nothing is asserted and left standing. Every formula in this book
is constructed in front of you --- usually starting from a version that is \emph{wrong},
watching it fail on a concrete example, and repairing exactly the part that failed. The
famous PUCT formula that governs Leela's every decision,
\[
  \argmax_a\;\Big[\,\underbrace{Q(s,a)}_{\text{what I've measured}}
  \;+\; \underbrace{c_{\text{puct}}\,P(a\given s)\,\frac{\sqrt{N(s)}}{1+N(s,a)}}_{\text{what I haven't checked yet}}\Big]
\]
is not introduced. It is \emph{assembled}, in Chapter~\ref{ch:puct}, out of five failed
attempts, each of which teaches you what one symbol is for. If you skip to the end and
read the formula, you have learned nothing. That is exactly the failure mode this book
exists to fix.

\section*{Who it is for}

You. A strong club player (2100--2200) who directs an engine-based training project, who
wants the theory properly rather than as a listicle, and who has been burned by
plausible-sounding explanations that turned out to be hollow.

\paragraph{Assumed:} chess to a strong club standard; secondary-school algebra (you can
rearrange $y = mx + c$ and you are not frightened by a summation sign).

\paragraph{Not assumed:} calculus, probability, linear algebra, machine learning, or
programming. Every mathematical tool used is built here, in the order it is needed. Where
calculus is genuinely required (once, for gradient descent, in Chapter~\ref{ch:learned}) it
is developed from the idea of a slope on a hill, with numbers.

\section*{The seven rules this book follows}

\begin{enumerate}
\item \textbf{Every formula is built, never dropped.} A formula arrives as version~0
  (crude, obviously wrong), is broken on an example, and is repaired. You see the scar
  tissue. That is the point: the final formula's strange parts --- why $\sqrt{N}$ and not
  $N$, why $1+N(s,a)$ and not $N(s,a)$ --- are exactly the repairs.

\item \textbf{Every idea gets arithmetic.} Not "the bonus decays"; instead, a table where
  the bonus is $0.55$, then $0.28$, then $0.19$, and you can check the division yourself.
  Understanding that survives contact with a real engine log is arithmetic-shaped.

\item \textbf{The engine numbers in this book are real.} Every table labelled
  \emph{Real engine output} was produced by running the actual \texttt{lc0.exe} and network
  in \texttt{engine/} on this machine, by the script
  \texttt{tools/collect\_engine\_data.py}. Nothing is typed from memory or imagination.
  Where a number is invented for the sake of clean arithmetic, it says so in the same
  breath. See "The honesty apparatus" below --- this matters more than it sounds.

\item \textbf{Chess claims are verified by an engine, not asserted.} This book teaches
  computer science, but it uses chess positions. Any chess claim ("this wins", "this
  throws away the win") has been checked with Stockfish at depth 30 and the check is
  reported. The author of the text is not a chess authority and does not pretend to be.

\item \textbf{Exercises are part of the exposition, not decoration.} Some results are
  \emph{only} developed in the exercises. Full solutions are in Appendix~\ref{app:solutions},
  worked at the same level of detail as the text. Do them with a pen. The arithmetic is
  deliberately small enough to do by hand.

\item \textbf{Confidence is labelled.} Textbook mathematics, published research findings,
  and this project's own untested hypotheses are three different things, and this book
  never lets them wear the same clothes.

\item \textbf{Everything lands somewhere in your repository.} Each chapter ends by pointing
  at the file in \texttt{chess\_speak\_out\_loud} where the idea actually lives, or notes
  honestly that it does not exist yet. Appendix~\ref{app:repomap} is the full map.
\end{enumerate}

\section*{The honesty apparatus}

This project has been burned twice by fluent, plausible, fabricated content --- a
"sacrifice" metric that never checked material, and a batch of invented master annotations.
The doctrine that came out of it (\emph{a bad coach does more harm than no coach}) applies
to a textbook at least as much as to a chess tutor. So claims in this book are colour-coded
by how much weight they will bear:

\begin{established}
Mathematics or documented engine behaviour. True by proof or by reading the source. If
this is wrong, the error is mine and it is a plain error.
\end{established}

\begin{researchfinding}
A published experimental result from the interpretability literature. Real evidence, real
peer scrutiny --- but obtained on a \emph{particular} network, with a particular method, and
not automatically true of the exact net in \texttt{engine/}. Reproduce before relying on it.
\end{researchfinding}

\begin{hypothesis}
An idea this project intends to build, which nobody has yet demonstrated works here. It
may be excellent. It is not knowledge. Where a chapter proposes one, it also proposes the
experiment that could kill it.
\end{hypothesis}

\begin{realdata}
Numbers produced by running the engine in \texttt{engine/} while writing this book. The
exact command line is in Appendix~\ref{app:enginedata}; re-run it and you should get the
same table (the search is deterministic at fixed node counts with one thread).
\end{realdata}

\noindent And, for the fourth case:

\begin{worked}{invented numbers, and why that is fine here}
When the point is the \emph{arithmetic of the algorithm}, real numbers are often clumsy ---
$P = 0.4413$ obscures the division you are meant to see. In those places this book uses
round invented numbers and says so in the caption: \emph{illustrative, not engine output}.
The rule is simple: invented numbers are allowed to demonstrate a \emph{mechanism}; they
are never allowed to support a \emph{claim about chess or about Leela}.
\end{worked}

\section*{How to read it}

The book has seven parts and a spine. The spine is Parts~\textsc{ii} and \textsc{iv}: the
search (how the engine decides where to spend its thinking) and the network (how it forms
an opinion in the first place). Everything else hangs off those.

\begin{center}
\begin{tikzpicture}[
   node distance=6mm,
   part/.style={draw, rounded corners=2pt, align=left, inner sep=5pt, font=\small,
                text width=0.72\textwidth},
   spine/.style={part, fill=BuildBlue!10, draw=BuildBlue},
   other/.style={part, fill=black!3, draw=black!35}]
\node[other] (p1) {\textbf{I. The problem} \hfill ch.\ 1--2\\
  \footnotesize what must be computed at all; why probability, not centipawns};
\node[spine, below=of p1] (p2) {\textbf{II. Searching by gambling well} \hfill ch.\ 3--7\\
  \footnotesize slot machines $\to$ UCB1 $\to$ UCT $\to$ PUCT $\to$ a full search done by hand};
\node[other, below=of p2] (p3) {\textbf{III. The search as it really runs} \hfill ch.\ 8--9\\
  \footnotesize FPU, batching, noise, tree reuse; how to read a real tree};
\node[spine, below=of p3] (p4) {\textbf{IV. The oracle} \hfill ch.\ 10--12\\
  \footnotesize hand-made formula $\to$ learned formula $\to$ attention $\to$ the two heads};
\node[other, below=of p4] (p5) {\textbf{V. Opening the skull} \hfill ch.\ 13--15\\
  \footnotesize probes, attention maps, learned look-ahead and prior override};
\node[other, below=of p5] (p6) {\textbf{VI. Bending the engine} \hfill ch.\ 16--17\\
  \footnotesize style as an objective function; sharpness that is real};
\node[other, below=of p6] (p7) {\textbf{VII. Making it speak} \hfill ch.\ 18--19\\
  \footnotesize divergence, facts, salience; one position end to end};
\draw[-{Latex[length=2mm]}, BuildBlue, thick] (p1) -- (p2);
\draw[-{Latex[length=2mm]}, BuildBlue, thick] (p2) -- (p3);
\draw[-{Latex[length=2mm]}, BuildBlue, thick] (p3) -- (p4);
\draw[-{Latex[length=2mm]}, BuildBlue, thick] (p4) -- (p5);
\draw[-{Latex[length=2mm]}, BuildBlue, thick] (p5) -- (p6);
\draw[-{Latex[length=2mm]}, BuildBlue, thick] (p6) -- (p7);
\end{tikzpicture}
\end{center}

\paragraph{Three ways through.}
\begin{itemize}
\item \textbf{Cover to cover} (the intended path). Parts build strictly on each other.
  Chapter~\ref{ch:byhand} --- a complete search executed by hand with real priors --- is
  the point where the first half either clicks or does not; do not pass it until it does.
\item \textbf{Search first, network later.} Chapters 1--9 are self-contained if you accept
  the network as a black box that returns two things. Many people find the search easier
  and the network more mysterious; this order lets the mystery wait.
\item \textbf{Reference.} Appendix~\ref{app:notation} is a one-page symbol table,
  Appendix~\ref{app:glossary} a glossary, Appendix~\ref{app:repomap} the map from concept
  to source file. Chapter~\ref{ch:readingtree} is the one to reread whenever you are
  staring at a confusing engine log.
\end{itemize}

\paragraph{Pace.} There are 19 chapters and roughly 130 exercises. At one chapter per
sitting this is about three weeks. The two long chapters (\ref{ch:byhand} and
\ref{ch:attention}) are each worth two sittings on their own.

\section*{A note on what this book will not do}

It will not do chess analysis. When a position appears, its role is to make an algorithm
concrete, and any evaluative statement about it is either (a) an engine output, quoted and
attributed, or (b) elementary and flagged as such. The governing rule of the project it
was written for --- \emph{the engine is the chess authority; language is a translator,
never a reasoner} --- is also the rule this book obeys about itself.
